\section*{Problemas}

\subsection*{Enunciados de los problemas}

% Recordar
\begin{problema}
	\label{prob:c7ea911}
	Una fábrica de baterías afirma que la duración promedio de sus baterías es de 500 horas con una desviación estándar poblacional conocida de $\sigma = 40$ horas. Un inspector de calidad toma una muestra aleatoria de $n = 64$ baterías y encuentra una duración promedio muestral de $\bar{X} = 490$ horas.
	\begin{enumerate}
		\item Formule las hipótesis nula ($H_0$) y alternativa ($H_a$) para probar si la duración promedio es estadísticamente menor a 500 horas.
		\item Calcule el estadístico de prueba $Z$.
		\item Determine el valor $p$ (p-valor) asociado a la prueba.
		\item ¿Se debe rechazar la hipótesis nula al nivel de significación $\alpha = 0.05$?
	\end{enumerate}
\end{problema}

% Comprender
\begin{problema}
	\label{prob:f6ad523}
	Un fabricante de neumáticos afirma que sus llantas duran en promedio $\mu_0 = 50,000$ km. Un consumidor toma una muestra aleatoria de $n = 16$ llantas y encuentra una media muestral $\bar{X} = 48,500$ km con una desviación estándar muestral $s = 3,000$ km.
	\begin{enumerate}
		\item ¿Qué distribución debe usar para la prueba de hipótesis ($Z$ o $t$ de Student)? Explique por qué.
		\item Formule las hipótesis para probar si la duración promedio poblacional es menor a 50,000 km.
		\item Calcule el estadístico $t$.
		\item Determine el valor $p$ aproximado usando $\nu = n-1 = 15$ grados de libertad.
		\item ¿Se debe rechazar $H_0$ al nivel de significación $\alpha = 0.05$?
	\end{enumerate}
\end{problema}

% Aplicar
\begin{problema}
	\label{prob:4072ece}
	Un fabricante de conductores eléctricos afirma que la resistencia promedio de sus cables industriales es exactamente de $\mu_0 = 100$ ohmios. Un ingeniero de pruebas toma una muestra de $n = 9$ tramos de cable y obtiene las siguientes lecturas de resistencia (en ohmios):
	\begin{center}
		98, 102, 99, 101, 97, 103, 100, 98, 102
	\end{center}
	\begin{enumerate}
		\item Calcule algebraicamente la media muestral $\bar{X}$ y la cuasivarianza muestral $s^2$.
		\item Formule las hipótesis para una prueba bilateral sobre la resistencia poblacional $\mu$.
		\item Calcule el estadístico $t$ de Student e indique sus grados de libertad.
		\item Determine el valor $p$ y concluya sobre la afirmación del fabricante al nivel $\alpha = 0.05$.
	\end{enumerate}
\end{problema}

% Analizar
\begin{problema}
	\label{prob:6c535e4}
	\textbf{Deducción teórica de invarianza de escala del estadístico $t$ de Student:} en la teoría de distribuciones exactas del muestreo normal, si $Z \sim N(0,1)$ y $V \sim \chi^2(\nu)$ son dos variables aleatorias independientes, el estadístico de Student se define como el cociente transformado:
	\begin{align}
		T = \frac{Z}{\sqrt{V / \nu}} \sim t_\nu.
	\end{align}
	Sea $X_1, \dots, X_n$ una muestra aleatoria normal $N(\mu, \sigma^2)$. Sabemos que $Z = \dfrac{\bar{X} - \mu}{\sigma/\sqrt{n}} \sim N(0,1)$ y que $V = \dfrac{(n-1)S^2}{\sigma^2} \sim \chi^2(n-1)$, siendo $\bar{X}$ y $S^2$ independientes por el Teorema de Cochran. Demuestra algebraicamente paso a paso que al sustituir estas expresiones en la definición de $T$, el parámetro de escala poblacional desconocido ($\sigma$) se cancela de manera exacta, demostrando formalmente por qué la variable pivotal $\displaystyle T = \frac{\bar{X} - \mu}{S/\sqrt{n}}$ permite realizar pruebas exactas sobre $\mu$ sin necesitar el conocimiento previo de $\sigma$.
\end{problema}

% Evaluar
\begin{problema}
	\label{prob:a9a27c5}
	Una compañía de telecomunicaciones asegura que el tiempo promedio de espera de un usuario en línea de atención al cliente es $\mu_0 = 8$ minutos con una desviación estándar $\sigma = 2$ minutos. Una firma auditora toma una muestra de $n = 49$ llamadas y encuentra un promedio muestral $\bar{X} = 8.5$ minutos.
	\begin{enumerate}
		\item Calcule el estadístico $Z$ para probar la hipótesis de que el tiempo de espera real excede los 8 minutos, y determine el valor $p$ de la prueba unilateral.
		\item ¿Cuál es la decisión y conclusión estadística al nivel de significación $\alpha = 0.05$?
		\item Evalúe si cambia la decisión cuando la auditoría exige un nivel de significación más estricto de $\alpha = 0.01$, y explique la implicación en términos del riesgo de error.
	\end{enumerate}
\end{problema}

% Crear
\begin{problema}
	\label{prob:8843fed}
	Diseñe un problema original de prueba de hipótesis sobre una media poblacional (distinto de baterías, llantas, cables o tiempos de espera ya vistos en esta sección): un laboratorio de control de calidad de alimentos afirma que el contenido de sodio promedio en un producto empacado es de $800$ mg. Un analista independiente toma una muestra de $n=36$ paquetes con contenido promedio muestral $\bar{X}=815$ mg, y desviación estándar poblacional conocida $\sigma=45$ mg. Especifique si debe usarse $Z$ o $t$, formule las hipótesis para una prueba bilateral, calcule el estadístico y el valor $p$, y concluya al nivel $\alpha=0.05$ si el contenido real difiere de $800$ mg.
\end{problema}

\subsection*{Soluciones detalladas}

% Recordar
\begin{solproblema}[prob:c7ea911]
	\begin{enumerate}
		\item Hipótesis unilaterales de cola izquierda: $H_0: \mu = 500$ horas, $H_a: \mu < 500$ horas.
		\item Estadístico $Z$ (dado que $\sigma=40$ es conocido):
		\begin{align}
			Z = \frac{\bar{X} - \mu_0}{\sigma / \sqrt{n}} = \frac{490 - 500}{40 / \sqrt{64}} = \frac{-10}{5} = -2.0.
		\end{align}
		\item $p\text{-valor} = P(Z < -2.0) = \Phi(-2.0) \approx 0.0228$ (2.28\%).
		\item Como $0.0228 < 0.05 = \alpha$, \textbf{rechazamos $H_0$}: existe evidencia suficiente para afirmar que la duración promedio es inferior a 500 horas.
	\end{enumerate}
\end{solproblema}

% Comprender
\begin{solproblema}[prob:f6ad523]
	\begin{enumerate}
		\item Debe usarse la \textbf{distribución $t$ de Student}: la desviación estándar poblacional $\sigma$ es desconocida (se estima con $s=3,000$ km) y el tamaño de muestra es pequeño ($n=16<30$), lo que impide invocar directamente el TLC con normalidad asintótica.
		\item $H_0: \mu = 50,000$ km, $H_a: \mu < 50,000$ km.
		\item $t = \dfrac{\bar{X}-\mu_0}{s/\sqrt{n}} = \dfrac{48,500-50,000}{3,000/\sqrt{16}} = \dfrac{-1,500}{750} = -2.0$.
		\item Con $\nu=15$: $p\text{-valor} = P(t_{15}<-2.0) \approx 0.032$ (3.2\%).
		\item Como $0.032 < 0.05$, \textbf{rechazamos $H_0$}: la vida útil real es inferior a la publicitada.
	\end{enumerate}
\end{solproblema}

% Aplicar
\begin{solproblema}[prob:4072ece]
	\begin{enumerate}
		\item $\sum X_i = 98+102+99+101+97+103+100+98+102=900$, $\bar{X}=900/9=100$ ohmios. Con $\sum(X_i-100)^2 = 4+4+1+1+9+9+0+4+4=36$: $s^2 = 36/8=4.5$, $s=\sqrt{4.5}\approx2.1213$ ohmios.
		\item $H_0: \mu=100$ ohmios, $H_a: \mu\neq100$ ohmios.
		\item $t = \dfrac{100-100}{2.1213/\sqrt9} = 0$, con $\nu=n-1=8$.
		\item $p\text{-valor} = 2\times P(t_8>0) = 2(0.5)=1.0$. Como $1.0>0.05$, \textbf{no rechazamos $H_0$}: la resistencia media coincide exactamente con la especificación de 100 ohmios.
	\end{enumerate}
\end{solproblema}

% Analizar
\begin{solproblema}[prob:6c535e4]
	Sustituyendo en el cociente $T = \dfrac{Z}{\sqrt{V/\nu}}$:
	\begin{align}
		T = \frac{\dfrac{\bar{X}-\mu}{\sigma/\sqrt{n}}}{\sqrt{\dfrac{(n-1)S^2/\sigma^2}{n-1}}}.
	\end{align}
	Simplificando el denominador: $\sqrt{\dfrac{(n-1)S^2}{\sigma^2(n-1)}} = \sqrt{S^2/\sigma^2} = S/\sigma$. Reescribiendo:
	\begin{align}
		T = \frac{\dfrac{\bar{X}-\mu}{\sigma/\sqrt{n}}}{S/\sigma} = \frac{\bar{X}-\mu}{\sigma/\sqrt{n}}\times\frac{\sigma}{S}.
	\end{align}
	El parámetro $\sigma$ se cancela aritméticamente entre numerador y denominador:
	\begin{align}
		T = \frac{\bar{X}-\mu}{S/\sqrt{n}} \sim t_{n-1}.
	\end{align}
	Esta cancelación exacta de $\sigma$ es lo que dota de poder a la inferencia de Student: permite que $T$ sea completamente calculable a partir exclusivamente de los datos de la muestra ($\bar{X}, S, n$) y de la hipótesis bajo prueba ($\mu_0$), bajo una distribución de referencia ($t_{n-1}$) 100\% independiente del verdadero valor de $\sigma$.
\end{solproblema}

% Evaluar
\begin{solproblema}[prob:a9a27c5]
	\begin{enumerate}
		\item $H_0:\mu=8$ vs. $H_a:\mu>8$. $Z = \dfrac{8.5-8.0}{2/\sqrt{49}} = \dfrac{0.5}{2/7} = 1.75$. $p\text{-valor} = P(Z>1.75) = 1-\Phi(1.75) = 1-0.9599 = 0.0401$ (4.01\%).
		\item Al nivel $\alpha=0.05$: como $0.0401<0.05$, \textbf{se rechaza $H_0$}: el tiempo de espera real excede los 8 minutos prometidos.
		\item Al nivel $\alpha=0.01$: como $0.0401>0.01$, \textbf{no se rechaza $H_0$}. Al exigir un riesgo de error Tipo I de solo el 1\%, la prueba requiere una divergencia muestral aún mayor para declarar rechazo ($Z>2.326$). Este caso ilustra cómo los resultados con valores $p$ intermedios (entre $0.01$ y $0.05$) son sensibles a la elección de $\alpha$: la misma evidencia muestral produce conclusiones opuestas según qué tan estricto sea el umbral de riesgo que la auditoría esté dispuesta a tolerar.
	\end{enumerate}
\end{solproblema}

% Crear
\begin{solproblema}[prob:8843fed]
	Como $\sigma=45$ mg es conocido, se usa el estadístico $Z$ (no $t$), independientemente del tamaño de muestra. $H_0:\mu=800$ mg, $H_a:\mu\neq800$ mg (bilateral).
	\begin{align}
		Z = \frac{\bar{X}-\mu_0}{\sigma/\sqrt{n}} = \frac{815-800}{45/\sqrt{36}} = \frac{15}{7.5} = 2.0.
	\end{align}
	Para una prueba bilateral:
	\begin{align}
		p\text{-valor} = 2\times P(Z>2.0) = 2\times[1-\Phi(2.0)] = 2(0.0228) = 0.0456.
	\end{align}
	Como $0.0456 < 0.05 = \alpha$, \textbf{se rechaza $H_0$}: existe evidencia estadística de que el contenido real de sodio difiere de los 800 mg declarados.
\end{solproblema}
